% 04 Apr cleanup


% Vocation and the Right to Interpret [editorial title]

\subsection{Installment 10: “Folkebladet” May 15, 1918}

It is remarkable that there are not a few in the Free Church who still adhere to neo-Lutheranism’s concept of the Church, as it was maintained by the once powerful theological faculty in Lund, Sweden—by such men as Professors Flensburg, Reuterdahl, Sundberg, Billing. In contrast to what Sverdrup taught in “Dogmatiske paragrafer” (S. S. VI), the Lund professors maintained that “the administration of the means of grace is not a product of the congregation’s life of faith, but rather its source; it is the Lord’s own institution. The Lord has not only ordered their function, but together with them an office for their administration, and the one who has not received this office by the ordered way has no right to administer the means of grace. If he does so, he breaks not only a human, but a divine order.”

Is it “the irony of fate” that this—the so-called “Lundensian church-concept,” over which there has been so much strife in Sweden, and which a portion of the Swedish Augustana Synod bequeathed to the Conference—still is upheld in the Free Church, whose principle no. 2 is to a great extent a negation of Lund theology’s concept of the Church?

Principle no. 2—that was the core of Birkeland’s articles and mine. These articles emphasized the religious (not formal) equality of the laity with the pastors, grounded in the principle of the universal priesthood. For most laypeople this was self-evident. But some pastors were thrown into confusion by it. The Board of Home Missions lost its head to such a degree that it connected a young pastor’s “fall” with these articles and with their authors’ influence upon the students at Augsburg. A semi-secret meeting of several pastors—both local and nonlocal—had been held in Minneapolis to go after the ordination articles and, naturally, their authors. The result was a declaration that was read before the meeting in Fargo. It served the reactionaries’ cause, was spiritless, and received the ambiguous epitaph, ‘It had already had its effect.’ A pity it was not allowed to stand in the annual report as a worthy and lasting piece of salt.\footnote{On Professor Helland’s proposal it had been decided that the declaration should be printed in the annual report. I supported the proposal (as did several others). It was adopted. At a later session two pastors proposed “that this decision be reconsidered, and that it be recommended to the signatories that they permit this to take place. We ground this on the fact that this resolution has already done its work...”

This proposal was not in order, insofar as it recommended to the signatories (Pastor Caspersen, etc.) that they grant permission for reconsideration. The declaration was, after all, the property of the meeting. It could dispose of the property as it wished, without asking any of the signatories’ advice. The comic thing was that a couple of the signatories, at the chairman’s prompting, rose and gave their “permission.” Did the chairman lose his bearings, or the movers, or the signatories—or all of them? A peculiar parliamentarian, very considerate besides.}

But to what end is it to contend over who it was that won? It is of very little significance to establish who the victors were. He who stands in accord with the truth is the victor, even if outwardly he suffers defeat. With that both parties may comfort themselves. For my own part I will say that no annual meeting decision can nullify the truth in these articles, which were both biblical and Lutheran.

These articles combated the neo-Lutheran, Romanizing concept of office, which the new direction had always combated. Among other things, they showed indirectly that the annual meetings in Brainerd and Marinette had fumbled in regard to ordination. And directly they reproduced the genuinely Lutheran doctrine, which has never been foreign to Augsburg.

For in 1895 Sverdrup wrote that ordination was not “a priestly matter,” not “an act of the gathered priesthood, but a Christian act, which can be performed by any congregation” (Veiledning 27); that it was “a congregational act, which only for reasons of expediency is carried out by several congregations in common” (25).

Evidently not very many at the annual meeting in Marinette were aware of these statements, where it was said that at Augsburg a new doctrine of ordination was being advanced. Evidently the ordainer at the annual meeting in Brainerd regarded ordination as a priestly matter, since he set aside two laymen (whom the meeting had elected to assist the ordainer, together with a couple of other pastors) and on his own authority appointed two pastors to function in their place. Evidently there is still confusion on the matter. A layman recently asked whether ordination would lose anything of its force if several laypeople were to take part in the laying on of hands. The answer he received was that Sverdrup, to be sure, was a layman, but he never took part in ordaining anyone. So then, what S. did or did not do was to determine ecclesiastical practice for generations to come!

With respect to the laying on of hands itself—well, it was not misconstrued in our articles. Our view of it was by no means foreign to Augsburg. As early as 1875 Sverdrup felt prompted to write the following, supported by Richard Rothe’s work, “Die Anfänge der chr. Kirche und ihrer Verfassung,” 160f. (S. has underlined important places in this work concerning polity):

\begin{quote}For the laying on of hands in the apostolic Church is not an ordination in our sense of the word, as clearly appears from Acts 8:17, where the laying on of hands accompanies the prayer over all the baptized, that they might receive the Holy Spirit—from Acts 13:3, where the congregation in Antioch prays with laying on of hands for Paul and Barnabas, who are sent out to the work of missions, after they have already long preached the word—and from 1 Tim. 4:14 and 2 Tim. 1:6, according to which it is at any rate highly probable that Timothy twice received laying on of hands without its ever being said that by this he received any power of office. In the same way we also understand the laying on of hands here, where there is mention of the overseers of the poor, as a special direction of the intercession toward the individual. (Samlede Skrifter I, 92).
\end{quote}

Even the great German stiffly orthodox Missouri Synod, on which the Free Church has not lavished many words of praise, is so free that it teaches:

\begin{quote}It is the call of the congregation that constitutes ministers, and actually confers the ministerial office. Ordination is not a divine ordinance, but an apostolic-ecclesiastical institution. It does not confer the ministry, as Papists and Romanizing Protestants assert, but it is only a public testimony and confirmation of the call. Ordination is not essential to the validity of the ministerial office.”
\end{quote}

I do not know what the editor teaches about ordination, but I think the annual meeting got the understanding that he would not agree that Augsburg had the view of ordination that B. and I had maintained.

3. It seems to me that Birkeland also has good reason to take a questioning and critical stance toward the editorial article “Church and the Churches,” an extract from the Lutheran’s reproduction of Dr. H. E. Jacobs’s article in “Constructive Quarterly.”

For many years I have known Jacobs’s concept of the Church. It is that of the General Council, concerning which in the East there has been strife for more than fifty years and of whose principle Oftedal wrote: it “is very far from Augsburg’s.” I was therefore painfully surprised when I saw Folkebladet’s editorial article “Church and the Churches,” and most of all when I read the editorial board’s assent to the quotations from Jacobs: “What is said here, we assent to with our whole heart in the Lutheran Free Church.”

No one who knows Jacobs can write like that, and who in more than a merely formal sense is in agreement with the Guiding Principles.

For with Dr. Jacobs doctrine is decisive in the church question. It is determinative for the boundaries of the confessional church and the particular church. The question of organization is for Jacobs merely and simply a question of expediency. The local congregation is to him no more divine than a church body; it is not the Kingdom’s proper form, for the Kingdom of God has many forms. Jacobs is also an eager spokesman for what we call “point 3.” He believes in the divine institution of the office, and finds the representative system in the New Testament. It was Jacobs’s view of the particular church that prevented the English-speaking Lutheran General Synod (1,400 pastors) from obtaining a place in the proposed general Lutheran conference in Philadelphia, because it (the Synod) did not accept the whole Book of Concord! I have written a couple of pamphlets in English about this. It takes a great deal of imagination to make Jacobs’s concept of the Church harmonize with the Free Church’s. And nevertheless—the editorial board wholeheartedly endorses the Lutheran’s extract from Jacobs’s article. I suppose the editorial board intended by its article “The Church and the Churches” to emphasize the whole Christian congregation—and that may be needed. But according to the Augsburg principle, the whole Christian congregation exists only in the form of individual congregations, not as individuals. And this emphasis on the whole Christian congregation becomes nothing other than what Dr. Schmidt emphasized, and which called forth S. and O.’s disapproval. In the light of this, and of what we have maintained concerning the principles, and of much else that we have not expressly cited—for example S. S. II, 222 and Veiledning 36–38—we see how difficult it is for the editorial board to steer by principle no. 1. We will see it even more if, for example, we should adopt the definition of the Church given by the late missionary pastor and professor Johs. Johnson.

In a lecture, 1915, Johnson says: “The Church is the common name for all the activity that the name of Jesus calls forth. From the most stiff-legged institution to the most unchurchly private endeavor to make God the center of human life.” An appealing definition—but very far from the one that “stands in harmony with Augsburg’s past.”

With these things before our eyes—and I could mention many others—it seems to me that Birkeland, who has been almost iron-consistent in his interpretation of the principles, has reason to ask, Is Folkebladet against Augsburg? By this question he has hardly meant to say that the one who thinks otherwise than the principles thinks unbiblically. But he has certainly meant to establish that only one authentic interpretation of the principles is possible, and that this interpretation in recent times has received headwind rather than fair wind in Folkebladet’s editorial articles.

By “Augsburg” B. must surely mean the old teachers who framed the principles. For now, it is evidently very different with the interpretation of the principles at Augsburg. Some think that principle no. 1 concerns the local congregation. Some teach that it concerns the holy catholic Church. And some perhaps take the golden middle way and think that it concerns both.

The same difference of opinion appears also among the pastors.

What applies to the different interpretations of principle no. 1 can also apply to the other principles. One solemnly professes them, and yet has different interpretations of them. But only one interpretation can be the right one. Now, according to the editorial board, the laity has “the right and call to interpret the Word of God,” and therefore they should have the right and indeed ought to have access to express themselves concerning the principles—in Folkebladet. The editorial board understands Birkeland’s interpretation and mine as dictatorship and absolutism—and wishes to hear from the laity. But:

The editorial board heard from the laity at the last annual meeting, but did not listen closely enough—so thought Fostervold.

With respect to the right and the call, it must be so that the right is conditioned by the call, and the call by the gift of grace. But is the gift of grace that enables one to interpret the Guiding Principles the common to all? Might it not be that those who have shown that they possess something of this gift are more qualified for this interpretation than those who have displayed only a superficial acquaintance with them?

If the principles are an intellectual deduction from the Word, then surely those who have thought most about them and have applied themselves most to understand them in the light of Scripture and church-historical development must be the more entitled to interpret them.

Now everyone does indeed have the right to have his opinion about the principles. Birkeland has it. The editorial board has it. A fourteen-year-old confirmand has it too. But I will not say that the fourteen-year-old has a call and therefore the right to write in Folkebladet about Guiding Principles as, for example, Birkeland and the editorial board do. Something does depend upon the regenerate personality—and that is true. But it depends likewise upon maturity, wide reading, and much else besides. Or has the twentieth century done away with the difference in gifts of grace?

With respect to B. and the editorial board, it seems to me that B. has furnished quite different and much more reassuring guarantees for the understanding of the principles than the editorial board has done. To be sure, the editor has been both professor and student at Augsburg for more than a generation. I mention this because the editorial board ascribes to seniority and locality a certain special value. But Birkeland, for his part, is not without redeeming qualifications.

Birkeland too, after all, has been at Augsburg—but likewise out in the wider world—both here in this country and in distant lands. He has kept his eyes open, reaped profit from his association with many church-historically significant personalities also outside Augsburg’s narrow walls. He has read much ecclesiastical literature and produced much that entitles him to be heard in the ecclesiastical sphere at least as much as any of his opponents. There is force and originality in what he writes. No compromise-making. No foggy speech. But straight-lined thoughts and a firm grasp. As regards the interpretation of the principles, he has kept himself surely and firmly and steadily turned in the right direction. If he has at times shot beyond the mark—and who does not do that?—he has nevertheless held the direction.

But now the editorial board has taken Birkeland to school. And me as well. And those poor “learned works” on polity, which are to yield place to the unschooled layman’s judgment! How this reminds one of that colleague of Luther who left his professorship in Wittenberg, took his Bible to the untaught peasants and asked them to interpret it for him, since he thought that theological learning was a hindrance to the Spirit! Flattering for the peasants, but nothing but the sheerest enthusiasm, in Luther’s view, for Luther, who was not dualistic in his conception of Christianity.

It is precisely for men like the editor that I have again and again referred to the literature of polity, history’s commentary on Scripture’s doctrine of the universal priesthood, the foundation of the Guiding Principles. There lies documented in this literature the thinking and experience of many generations, from which the editorial board somewhat grandly distances itself, and of which his writing betrays so little knowledge.

There are many who understand the Guiding Principles without having specialized in polity literature. Laypeople especially, Norwegian academics as well. But how difficult it proves for a portion of the Augsburg men to understand them! Is it because Augsburg has not been able to give them enough light, or does the fault lie in themselves?
